\documentclass[book.tex]{subfiles}
\begin{document}
\chapter{Equivariant Dynamical Models}

\label{chap:eqdynamic}
\noindent\rule{11cm}{0.4pt} \\
\textbf{Prerequisites:} \autoref{chap:qft_basics}  \\
\textbf{Difficulty Level:} ** \\
\noindent\rule{11cm}{0.4pt}
\vspace{5mm}

Equivariances can be incorporated into models of dynamical systems in order to improve data efficiency \cite{wang2020incorporating}. A differential equation is said to respect a symmetry group $G$ if all solutions to the equation respect group $G$.

\section{Navier Stokes Symmetries}

Solutions to Navier-Stokes respect a number of symmetries
\begin{itemize}
\item Space translation
\item Time translation
\item Uniform shifts
\begin{align}
w(x, t) \mapsto \lambda w(x, t) + c
\end{align}
\item Rotation/Reflection
\item Scaling
\begin{align}
w(x, t) \mapsto \lambda w(\lambda x, \lambda^2 t)
\end{align}
\end{itemize}

We can use equivariant convolutions that are equivariant to these known symmetries when modeling fluid systems. 
%"""
%Recent work has shown deep learning can accelerate the prediction of physical dynamics relative to numerical solvers. However, limited physical accuracy and an inability to generalize under distributional shifts limit its applicability to the real world. We propose to improve accuracy and generalization by incorporating symmetries into convolutional neural networks. Specifically, we employ a variety of methods each tailored to enforce a different symmetry. Our models are both theoretically and experimentally robust to distributional shifts by symmetry group transformations and enjoy favorable sample complexity. We demonstrate the advantage of our approach on a variety of physical dynamics including Rayleigh–Bénard convection and real-world ocean currents and temperatures. Compared with image or text applications, our work is a significant step towards applying equivariant neural networks to high-dimensional systems with complex dynamics.
%"""



\end{document}